\documentclass[10pt,a4paper]{amsart}
\usepackage[margin=19mm]{geometry}
\usepackage{amsmath,amssymb,mathtools}
\usepackage[scr=rsfs,cal=euler]{mathalfa}
\usepackage{xfrac}
\usepackage[unicode,hidelinks,bookmarks=false]{hyperref}
\newcommand{\ie}{i.e.\ }
\newcommand{\cf}{cf.\ }
\newcommand{\resp}{respectively\ }
\newcommand{\Subsection}{\subsection}
\providecommand{\nobreakdash}{\nobreak\hbox{-}}
\setlength{\emergencystretch}{2em}
\multlinegap0pt
\usepackage{amssymb}
\usepackage{mathtools}
\usepackage{tikz}
\newtheorem{prop}{Proposition}
\newtheorem{lemma}{Lemma}
\newcommand{\Cs}{\mathscr{C}}
\newcommand{\Fs}{\mathscr{F}}
\newcommand{\Gs}{\mathscr{G}}
\newcommand{\IK}{\Is_K}
\newcommand{\Is}{\mathscr{I}}
\newcommand\N{\mathfrak{N}}
\newcommand\Nt{\widetilde{N}}
\newcommand{\OK}{\Os_K}
\newcommand{\Os}{\mathcal{O}}
\newcommand\ab{\boldsymbol{a}}
\newcommand\afr{{\mathfrak{a}}}
\newcommand\afrb{{\underline{\mathfrak{a}}}}
\newcommand\cfrb{{\underline{\mathfrak{c}}}}
\newcommand\sums[1]{\sum_{\substack{#1}}}
\DeclareMathOperator{\vol}{vol}
\title{Integral points over number fields: a Clemens complex jigsaw puzzle}
\author{Christian Bernert, Ulrich Derenthal, Judith Ortmann, Florian Wilsch}
\date{Source: arXiv:2601.08774v2 (CC BY 4.0)}
\begin{document}
\maketitle

\noindent\emph{Source and licence.} Integral points over number fields: a Clemens complex jigsaw puzzle. Version arXiv:2601.08774v2 (CC BY 4.0), distributed by arXiv under the Creative Commons Attribution 4.0 International licence, http://creativecommons.org/licenses/by/4.0/, as recorded in the arXiv licence metadata for this exact version and shown on the abstract page https://arxiv.org/abs/2601.08774v2. Full licence notice: \texttt{licenses/arxiv-2601.08774-CC-BY-4.0.txt}.\par\smallskip

\noindent\textit{Editorial scope.} Counted unit: the article's prop prop:elements\_to\_ideals (source0.tex lines 1081--1089). The article's complete original proof of it is delivered with it (source0.tex lines 1091--1101, one \texttt{proof} environment of 11 lines). No mathematics is written, repaired or re-derived: the statement and the proof are the article's own lines, in the article's own notation, and no retained line is substituted or re-flowed.  Retained uncounted as the source-local context the proof needs: lines 865--867; lines 869--871; lines 873--875; lines 939--941; lines 964--966; lines 968--970; lines 976--982; lines 1009--1014; lines 1044--1049.  Nothing was omitted from any retained span, so the package carries no omission record.  No retained line contains a citation, so the wrapper carries no bibliography and no bibliography entry.  No retained line contains a figure environment, an image-inclusion command or drawing code, so no image is delivered and the package declares no asset dependency.  Editorial formatting only, in addition: an AMS \texttt{amsart} preamble that restates the article's own declarations and macros used by the retained lines, namely the environments \texttt{prop}, \texttt{lemma} on separate counters, together with the standard packages those lines need.  The retained environments print in this document as Proposition 1, Lemma 1 in order of appearance, whereas the article prints the same items with the numbering of its own full document.  The complete native source of the article as distributed in the arXiv e-print for this exact version is bundled unchanged as \texttt{sources/192624/source0.tex}.
    \begin{equation}\label{eq:torsor}
            a_2a_8+a_1a_9+a_3a_4^2a_5^3a_7 = 0,
    \end{equation}

    \begin{equation}\label{eq:height}
        \prod_{v \mid \infty} \Nt_v(a_1^{(v)},\dots,a_9^{(v)})\le B,
    \end{equation}

    \begin{equation}\label{eq:fundamental_domain}
        (a_1,\dots,a_9) \in \Fs \subset (K^\times)^7 \times K^2,
    \end{equation}

\begin{equation}\label{eq:lattice_points}
    |\{(a_8,a_9) \in \Os_8 \times \Os_9 : \eqref{eq:torsor}, \eqref{eq:height}, \eqref{eq:fundamental_domain}\}| = |\Gs(\ab',a_6,a_7,\cfrb) \cap \Fs_0(\ab',a_6,a_7;B)|.
\end{equation}

\begin{equation}\label{eq:restrict_a6_a7}
    |a_i|_v \le (B_i/T)^{d_v/d}
\end{equation}

\begin{equation}\label{eq:restrict_a1a2}
    \N(\afr_1\afr_2) \le \frac{B}{T}
\end{equation}

\begin{prop}\label{prop:restrict_a6_a7}
    We have
    \begin{align*}
        |M_\cfrb(B)| ={}&\sums{\ab' \in \Os'_* \cap \Fs_1^5\\ \eqref{eq:restrict_a1a2}} \sums{a_6,a_7 \in \OK^\times\\\eqref{eq:restrict_a6_a7}} |\{(a_8,a_9) \in \Os_8 \times \Os_9 : \eqref{eq:torsor}, \eqref{eq:height}, \eqref{eq:fundamental_domain}\}|\\
        &+ O(B(\log B)^{1+2q}\log\log B).
    \end{align*}
\end{prop}

\begin{lemma}\label{lem:lattice_points}
    For $\ab' \in \Os'_*$ with $\theta_0(\afrb')=1$ and $a_6,a_7 \in \OK^\times$, we have
    \begin{align*}
        |\Gs(\ab',a_6,a_7,\cfrb) \cap \Fs_0(\ab',a_6,a_7;B)| ={}& \frac{2^{r_2}\vol S_F(\ab',(a_6^{(v)},a_7^{(v)})_v;B)}{|\Delta_K|^{1/2}\N(\afr_1\Os_8)}\\ &+ O\left(\frac{B}{T^{1/d}\N(\afr_1\afr_2)}\right),
    \end{align*}
\end{lemma}

\begin{lemma}\label{lem:vol_S_F}
    For $\ab' \in \Os_*' \cap \Fs_1^5$ with $\theta_0(\afrb')=1$ and $a_6,a_7 \in \OK^\times$ with \eqref{eq:restrict_a6_a7}, we have
    \begin{equation*}
        \vol S_F(\ab',(a_6^{(v)},a_7^{(v)})_v;B) = \frac{2^{r_1}\pi^{r_2}R_K B}{|N(a_2)|} \left(1 + O(T^{-1/d})\right).
    \end{equation*}
\end{lemma}

\begin{prop}\label{prop:elements_to_ideals}
    We have
    \begin{equation*}
        \begin{aligned}
            \sum_{\cfrb \in \Cs} |M_\cfrb(B)| ={}& \frac{2^{r_1}(2\pi)^{r_2}(1 + O(T^{-1/d})) R_K}{|\Delta_K|^{1/2}} \sums{(\afr_1,\afr_2) \in \IK^2\\\eqref{eq:restrict_a1a2}\\        [\afr_1\afr_2]=[\OK]} \frac{\theta_0(\afrb')B}{\N(\afr_1\afr_2)} \sums{a_6,a_7 \in \OK^\times\\\eqref{eq:restrict_a6_a7}} 1\\
        &+ O(B(\log B)^{1+2q}(\log\log B))
        \end{aligned}
    \end{equation*}
\end{prop}

\begin{proof}
    We combine Proposition~\ref{prop:restrict_a6_a7} with \eqref{eq:lattice_points}, Lemma~\ref{lem:lattice_points}, and Lemma~\ref{lem:vol_S_F}. Recall that $\ab' \in \Os_*'\cap \Fs_1^5$ implies $a_3=a_4=a_5=1$.
    
    Summing the error term from Lemma~\ref{lem:lattice_points} over the remaining variables $a_1,\dots,a_7$ yields a total contribution bounded by
    \begin{equation*}
        \ll \frac{1}{T^{1/d}} \sums{\ab' \in \Os'_* \cap \Fs_1^5\\\eqref{eq:restrict_a1a2}} \sums{a_6,a_7 \in \OK^\times\\ \eqref{eq:restrict_a6_a7}} \frac{B}{\N(\afr_1\afr_2)} \ll \frac{B(\log B)^{2+2q}}{T^{1/d}}
    \end{equation*}
    by an argument similar to the proof of Proposition~\ref{prop:restrict_a6_a7}, which is satisfactory.
   
    In the main term, we replace the summations over $(a_1,a_2) \in (\Os_{1*}\times\Os_{2*})\cap \Fs_1^2$ and $\cfrb \in \Cs$ by a summation over ideals $\afr_1,\afr_2$ such that $\afr_1\afr_2$ is a principal ideal (since $\Os_1\Os_2 = \OK$), and we use that $|N(a_2)|\N(\Os_8) = \N(\afr_2)$ (since $\Os_8 = \Os_2^{-1}$).
\end{proof}

\end{document}
